% =============================================================================
% T0_preamble_book_plus.tex — Ergänzung der Buch-Preamble für Bd. 6–8
% Die neueren Einzeldokumente (ab Dok. 263) wurden für die A4-Standalone-
% Preamble geschrieben. Diese Datei stellt die dort zusätzlich definierten
% Elemente auch im Buchsatz bereit und kapselt dokumenteigene Titel-,
% Inhaltsverzeichnis- und Kapitelbefehle, damit jedes Dokument genau ein
% Buchkapitel mit genau einem Eintrag im Inhaltsverzeichnis ergibt.
% =============================================================================

% =============================================================================
% T0_preamble_book_appendixfix.tex — \appendix innerhalb eines Buchkapitels
% Viele Einzeldokumente rufen \appendix auf. Im Buch darf das nicht auf die
% Kapitelebene durchschlagen (sonst heißen alle folgenden Dokumente
% "Anhang n"). Hier bewirkt \appendix nur, dass die Abschnitte des laufenden
% Dokuments mit Buchstaben gezählt werden; das nächste Kapitel setzt die
% normale Abschnittszählung wieder ein.
% =============================================================================
\makeatletter
\let\FFGFT@origthesection\thesection
\renewcommand\appendix{\par
  \setcounter{section}{0}%
  \gdef\thesection{\thechapter.\Alph{section}}}
\pretocmd{\chapter}{\global\let\thesection\FFGFT@origthesection}{}{}
\makeatother


% Bildpfade der neueren Dokumente
\graphicspath{{../Figuren/}{../wr_standalone_A4/}{../fig/}{../figures/}}

% Farb-Aliase aus der Standalone-Preamble
\providecolor{t0gray}{gray}{0.45}
\colorlet{T0blue}{t0blue}
\colorlet{T0red}{t0red}
\colorlet{T0green}{t0green}
\colorlet{T0orange}{t0orange}
\colorlet{T0purple}{t0purple}
\colorlet{T0yellow}{t0yellow}
\colorlet{T0gray}{boxgray}

% Seitenstil, den einige Dokumente aufrufen
\fancypagestyle{firstpage}{%
  \fancyhf{}%
  \fancyfoot[C]{\thepage}%
  \renewcommand{\headrulewidth}{0pt}%
}

% Breite Tabellen auf Satzspiegel begrenzen (wie A4-Standalone)
\BeforeBeginEnvironment{tabular}{\begin{adjustbox}{max width=\linewidth,center}}
\AfterEndEnvironment{tabular}{\end{adjustbox}}

\sloppy
\emergencystretch=4em

% Mathematik in Kapiteltiteln: lesbare PDF-Lesezeichen
\pdfstringdefDisableCommands{%
  \def\varphi{phi}\def\mitvarphi{phi}\def\Phi{Phi}\def\mitPhi{Phi}%
  \def\leftrightarrow{<->}\def\Leftrightarrow{<=>}%
  \def\hat#1{#1}\def\tilde#1{#1}\def\widetilde#1{#1}%
  \def\mathbb#1{#1}\def\mathrm#1{#1}\def\mathcal#1{#1}\def\mathbf#1{#1}%
  \def\text#1{#1}\def\operatorname#1{#1}%
  \def\sqrt#1{sqrt(#1)}\def\frac#1#2{#1/#2}\def\tfrac#1#2{#1/#2}%
  \def\textsuperscript#1{#1}\def\textsubscript#1{#1}%
  \def\infty{inf}\def\to{->}\def\mitzeta{zeta}\def\zeta{zeta}%
  \def\mitiota{iota}\def\iota{iota}\def\mitxi{xi}\def\mittheta{theta}%
}

% -----------------------------------------------------------------------------
% \FFGFTdoc{<Kapiteltitel>}{<ch-Datei>}      — normales Einzeldokument
% \FFGFTdocDemote{<Kapiteltitel>}{<ch-Datei>} — Dokument, das intern
%      \chapter als oberste Gliederungsebene benutzt (Dok. 308–314, 335):
%      \chapter -> \section, \section -> \subsection, \subsection -> \subsubsection
% -----------------------------------------------------------------------------
\makeatletter
\let\FFGFT@origacl\addcontentsline
\let\FFGFT@origchapter\chapter
\let\FFGFT@origsection\section
\let\FFGFT@origsubsection\subsection
\let\FFGFT@origsubsubsection\subsubsection
\def\FFGFT@chapterword{chapter}
% Mehrere Dokumente definieren dieselben Theorem-Umgebungen selbst;
% im Buch wird eine bereits vorhandene Umgebung stillschweigend neu angelegt.
\let\FFGFT@orignewtheorem\newtheorem
\def\FFGFT@nt#1#2{%
  \expandafter\let\csname #2\endcsname\relax
  \expandafter\let\csname end#2\endcsname\relax
  \expandafter\let\csname c@#2\endcsname\relax
  \FFGFT@orignewtheorem#1{#2}}
\renewcommand\newtheorem{\@ifstar{\FFGFT@nt*}{\FFGFT@nt{}}}
\let\FFGFT@origend\end
\def\FFGFT@docword{document}
\def\FFGFT@end#1{%
  \def\FFGFT@tmp{#1}%
  \ifx\FFGFT@tmp\FFGFT@docword
    \endinput\expandafter\@gobble
  \else
    \expandafter\FFGFT@origend
  \fi{#1}}
\newcommand\FFGFT@neutralize{%
  \let\maketitle\relax
  \let\tableofcontents\relax
  \let\listoftables\relax
  \let\listoffigures\relax
  \renewcommand\addcontentsline[3]{%
    \def\FFGFT@tmp{##2}%
    \ifx\FFGFT@tmp\FFGFT@chapterword\else\FFGFT@origacl{##1}{##2}{##3}\fi}%
}
\newcommand\FFGFTdoc[2]{%
  \FFGFT@origchapter{#1}%
  \begingroup
    \FFGFT@neutralize
    \renewcommand\chapter{\@ifstar{\@gobble}{\FFGFT@origsection}}%
    \input{#2}%
  \endgroup
}
\newcommand\FFGFTdocDemote[2]{%
  \FFGFT@origchapter{#1}%
  \begingroup
    \FFGFT@neutralize
    \renewcommand\chapter{\@ifstar{\@gobble}{\FFGFT@origsection}}%
    \let\section\FFGFT@origsubsection
    \let\subsection\FFGFT@origsubsubsection
    % Dok. 311–314 enden mit \end{document}: dort nur die Datei beenden
    \let\end\FFGFT@end
    \input{#2}%
  \endgroup
}
\makeatother
